\documentclass[11pt]{report}

\usepackage[margin=1in]{geometry}
\usepackage{amsmath}
\usepackage{graphicx}
\usepackage{hyperref}
\usepackage{algorithm}
\usepackage{algpseudocode}
\usepackage{float}
\title{Computational Analysis of SEEG Neural Population Dynamics}
\author{}
\date{}

\begin{document}

\maketitle

\tableofcontents

\chapter{Introduction}

Stereo-electroencephalography (SEEG) provides a unique opportunity to measure neural activity
directly from distributed cortical networks with high temporal resolution.
However, the high dimensionality and nonstationary nature of SEEG signals require
integrated computational frameworks that combine signal processing,
statistical inference, and population-level dynamical analysis.

This report presents a comprehensive computational framework for SEEG analysis
that integrates classical electrophysiological analysis methods with modern
population dynamics approaches. The framework includes preprocessing,
time--frequency analysis, statistical inference, latent trajectory analysis,
and dynamical systems modeling of multichannel neural activity.

\chapter{Experimental Data and Recording Methods}

This chapter describes the experimental setup and recording procedures
used to acquire stereo-electroencephalography (SEEG) signals.
We summarize electrode implantation strategies, recording hardware,
signal acquisition parameters, and experimental paradigms used to
evoke neural responses. In addition, we describe the procedures used
to synchronize experimental events with SEEG recordings.

\section{Participants}

SEEG recordings were obtained from participants undergoing clinical
evaluation for drug-resistant epilepsy. Electrode implantation was
performed based on clinical requirements determined by the treating
neurologists and neurosurgeons. All experimental procedures were
conducted in accordance with institutional ethical guidelines and
approved by the relevant ethics committees.

Participants provided informed consent for the use of their SEEG data
for research purposes.

\section{SEEG Electrode Implantation}

Stereo-electroencephalography electrodes were implanted using a
stereotactic surgical procedure. Each electrode consists of a linear
array of cylindrical contacts that sample neural activity along the
trajectory of the electrode.

Typical electrode characteristics include:

\begin{itemize}
\item Contact diameter: approximately 0.8--1.0 mm
\item Contact length: 2 mm
\item Inter-contact spacing: 3--5 mm
\item Number of contacts per electrode: 8--16
\end{itemize}

Electrodes were inserted to target cortical and subcortical regions
of interest based on clinical hypotheses regarding seizure onset zones.
The spatial distribution of electrode contacts provides sparse but
high-resolution sampling of neural activity across distributed brain
regions.

\section{Recording Hardware and Signal Acquisition}

SEEG signals were recorded using a clinical intracranial EEG
recording system. Neural signals were amplified, digitized, and stored
for offline analysis.

Typical acquisition parameters include:

\begin{itemize}
\item Sampling rate: 1000--2000 Hz
\item Analog bandwidth: approximately 0.1--500 Hz
\item Digital resolution: 16-bit or higher
\end{itemize}

Signals were initially referenced to a common clinical reference
electrode. Re-referencing was performed offline during preprocessing
to obtain bipolar or common-average referenced signals.

\section{Experimental Paradigms}

Two types of experimental paradigms were analyzed in this study:

\begin{itemize}
\item Electrical stimulation experiments
\item Cognitive task experiments
\end{itemize}

These paradigms provide complementary approaches to probing neural
dynamics in SEEG recordings.

\subsection{Electrical Stimulation Experiments}

Electrical stimulation experiments were performed using pairs of
adjacent SEEG contacts. Biphasic electrical pulses were delivered
to stimulate cortical tissue while neural responses were recorded
from surrounding electrodes.

Typical stimulation parameters include:

\begin{itemize}
\item Pulse shape: biphasic square wave
\item Pulse duration: 0.2--1 ms per phase
\item Current amplitude: 1--5 mA
\item Stimulation frequency: single pulse or low-frequency trains
\end{itemize}

Electrical stimulation provides controlled perturbations of neural
circuits, allowing characterization of evoked responses and network
propagation patterns.

\subsection{Cognitive Task Experiments}

In addition to electrical stimulation experiments, SEEG recordings
were obtained during cognitive task paradigms designed to probe
sensory, motor, or cognitive processing.

Tasks were implemented using stimulus presentation software that
recorded precise event timestamps for each trial. Typical events
include stimulus onset, participant responses, and task-specific
markers.

These event markers allow alignment of neural activity to task
events for subsequent analysis.

\section{Event Timing and Synchronization}

Precise synchronization between neural signals and experimental
events is essential for event-related analysis.

Event markers were transmitted from the stimulus presentation
computer to the SEEG recording system using hardware triggers.
These triggers were recorded alongside neural signals and stored
as digital event channels.

Each event entry contains the following information:

\begin{itemize}
\item Event type
\item Event timestamp
\item Trial identifier
\item Experimental condition label
\end{itemize}

Event timestamps were used to construct event-locked neural epochs
for subsequent analysis.

\section{Data Organization and Metadata}

SEEG datasets were organized into structured formats containing
neural signals, event markers, and channel metadata.

Typical dataset components include:

\begin{itemize}
\item Continuous SEEG signals
\item Channel labels
\item Sampling rate information
\item Event tables
\item Anatomical labels for electrode contacts
\end{itemize}

This structured organization facilitates reproducible computational
analysis and integration with signal processing pipelines.


\chapter{Signal Preprocessing}

Raw stereo-electroencephalography (SEEG) recordings contain neural
signals mixed with various sources of noise and artifacts.
Preprocessing is therefore required to remove non-neural components
while preserving physiologically meaningful neural activity.

The preprocessing pipeline consists of the following stages:

\begin{enumerate}
\item Filtering and signal conditioning
\item Re-referencing
\item Bad channel detection
\item Continuous artifact detection
\item Electrical stimulation artifact suppression
\end{enumerate}

These steps produce cleaned continuous SEEG signals that can be used
for event-related and population-level analyses.

\section{SEEG Signal Model}

The recorded SEEG signal can be modeled as a mixture of neural activity,
artifacts, and measurement noise:

\begin{equation}
x(t) = s(t) + a(t) + \epsilon(t)
\end{equation}

where

\begin{itemize}
\item $s(t)$ represents neural population activity
\item $a(t)$ represents artifacts (e.g., stimulation, motion, line noise)
\item $\epsilon(t)$ represents measurement noise
\end{itemize}

The goal of preprocessing is to estimate a cleaned signal

\begin{equation}
\tilde{x}(t) \approx s(t)
\end{equation}

while minimizing distortion of neural dynamics.

\section{Filtering}

Filtering removes slow drift and high-frequency noise from SEEG signals.

\subsection{High-Pass Filtering}

Low-frequency drift can arise from electrode polarization or slow
physiological fluctuations. A high-pass filter removes these slow
components:

\begin{equation}
y(t) = x(t) * h_{HP}(t)
\end{equation}

where $h_{HP}(t)$ is the impulse response of the high-pass filter.

Typical cutoff frequency:

\begin{itemize}
\item $f_{HP} = 0.5$ Hz
\end{itemize}

\subsection{Low-Pass Filtering}

High-frequency noise above the physiological bandwidth is suppressed
using a low-pass filter:

\begin{equation}
y(t) = x(t) * h_{LP}(t)
\end{equation}

Typical cutoff frequency:

\begin{itemize}
\item $f_{LP} = 250$ Hz
\end{itemize}

\subsection{Line Noise Removal}

Electrical line noise introduces oscillatory contamination at
50 Hz (or 60 Hz depending on location) and its harmonics.

A notch filter is applied at each harmonic frequency:

\begin{equation}
y(t) = x(t) * h_{notch}(t)
\end{equation}

Typical notch frequencies include:

\begin{itemize}
\item 50 Hz
\item 100 Hz
\item 150 Hz
\end{itemize}

\section{Re-Referencing}

SEEG recordings are initially referenced to a common clinical reference.
To obtain physiologically meaningful signals, offline re-referencing
is applied.

\subsection{Common Average Reference}

In common-average referencing (CAR), the average signal across all
channels is subtracted from each channel:

\begin{equation}
x_i^{CAR}(t) =
x_i(t) -
\frac{1}{N}\sum_{j=1}^{N} x_j(t)
\end{equation}

where $N$ is the number of channels.

\subsection{Bipolar Referencing}

Bipolar referencing computes the difference between adjacent electrode
contacts:

\begin{equation}
x_i^{bip}(t) = x_i(t) - x_{i+1}(t)
\end{equation}

Bipolar referencing suppresses common noise sources and improves
spatial localization of neural signals.

\section{Bad Channel Detection}

Some SEEG channels may contain excessive noise due to poor electrode
contact or hardware issues.

Channels are identified as bad if they exhibit abnormal statistical
properties.

Typical metrics include:

\begin{itemize}
\item Signal variance
\item RMS amplitude
\item Line noise ratio
\item Kurtosis
\end{itemize}

For example, channel variance is computed as

\begin{equation}
\sigma_i^2 = \frac{1}{T} \sum_{t=1}^{T} (x_i(t) - \bar{x}_i)^2
\end{equation}

Channels with variance exceeding a threshold relative to the median
variance across channels are excluded from further analysis.

\section{Continuous Artifact Detection}

Large-amplitude artifacts can contaminate SEEG signals during recording.

Artifact segments are detected using amplitude thresholds:

\begin{equation}
|x(t)| > \theta
\end{equation}

where $\theta$ is a predefined amplitude threshold.

Detected artifact segments can either be:

\begin{itemize}
\item removed from analysis
\item replaced using interpolation
\item masked during event-locked analysis
\end{itemize}

\section{Electrical Stimulation Artifact Suppression}

Electrical stimulation produces large transient artifacts that
dominate SEEG signals immediately following the stimulation pulse.

To suppress stimulation artifacts, a short post-stimulation time
window is masked:

\begin{equation}
t_{mask} = [t_{stim}, t_{stim} + \Delta t]
\end{equation}

where $\Delta t$ typically ranges from 5--10 ms.

Signal samples within this window are excluded from time–frequency
analysis and statistical tests.

\section{Output of Preprocessing}

The preprocessing pipeline produces the following outputs:

\begin{itemize}
\item Cleaned continuous SEEG signals
\item List of valid channels
\item Artifact masks
\item Metadata describing preprocessing parameters
\end{itemize}

These cleaned signals are subsequently used to construct event-locked
epochs for trial-based neural analysis.





\chapter{Event-Locked Signal Extraction}

Event-related analysis requires segmentation of continuous SEEG signals
into trials aligned to experimentally defined events. These events may
correspond to electrical stimulation pulses or behavioral task events.
Event-locked segmentation transforms continuous neural recordings into
a structured trial representation suitable for statistical and
time--frequency analysis.

Let $x_i(t)$ denote the continuous SEEG signal recorded from channel $i$.
Let $\tau_k$ denote the timestamp of event $k$.
Event-locked neural signals are obtained by extracting time windows
around each event.

\begin{equation}
x_i^{(k)}(t) = x_i(\tau_k + t)
\end{equation}

where $t$ denotes time relative to the event.

The resulting dataset consists of a three-dimensional tensor:

\begin{equation}
X(k,i,t)
\end{equation}

where

\begin{itemize}
\item $k$ indexes trials
\item $i$ indexes recording channels
\item $t$ indexes time samples relative to the event
\end{itemize}

This representation forms the basis for all subsequent analyses.

\section{Epoch Construction}

For each event timestamp $\tau_k$, an epoch is extracted from the
continuous signal using a predefined time window:

\begin{equation}
t \in [-T_{pre}, T_{post}]
\end{equation}

The extracted epoch therefore contains neural activity before and
after the event:

\begin{equation}
x_i^{(k)}(t) =
x_i(\tau_k + t),
\quad
t \in [-T_{pre}, T_{post}]
\end{equation}

Typical values include

\begin{itemize}
\item $T_{pre} = 1.0$ s
\item $T_{post} = 1.0$ s
\end{itemize}

Epoch extraction produces a dataset of size

\begin{equation}
K \times N \times T
\end{equation}

where

\begin{itemize}
\item $K$ is the number of trials
\item $N$ is the number of channels
\item $T$ is the number of time samples
\end{itemize}

\section{Trial Structure}

Each extracted trial contains three conceptual segments:

\begin{enumerate}
\item Prestimulus baseline
\item Event response period
\item Late post-event activity
\end{enumerate}

This structure allows separation of baseline neural activity from
event-driven responses.

\subsection{Prestimulus Window}

The prestimulus window serves as a baseline reference period.
Let the baseline interval be defined as

\begin{equation}
t \in [-T_{base}, 0]
\end{equation}

Baseline activity is used for:

\begin{itemize}
\item baseline normalization of spectral power
\item trial-level quality control
\item estimation of spontaneous neural activity
\end{itemize}

Typical baseline duration is

\begin{equation}
T_{base} = 500\text{--}800 \text{ ms}
\end{equation}

\subsection{Poststimulus Window}

The poststimulus interval contains the neural response to the event:

\begin{equation}
t \in [0, T_{post}]
\end{equation}

For electrical stimulation experiments, the earliest portion of this
window may contain stimulation artifacts. These samples are masked
before further analysis:

\begin{equation}
t \in [0, t_{mask}]
\end{equation}

where $t_{mask}$ typically ranges from 5--10 ms.

The remaining poststimulus interval contains physiologically meaningful
neural responses used for ERP and time--frequency analyses.

\chapter{Trial-Level Quality Control}

Trial-level quality control procedures identify trials contaminated by
noise or artifacts. Because SEEG recordings may contain transient
disturbances, quality control ensures that downstream analyses are
performed using reliable neural data.

Quality metrics are computed from prestimulus neural activity to
avoid bias from event-related responses.

\section{Prestimulus Signal Metrics}

Let $x_i^{(k)}(t)$ denote the signal from channel $i$ during trial $k$.
Metrics are computed within the baseline window:

\begin{equation}
t \in [-T_{base},0]
\end{equation}

These statistics characterize the stability and noise level of
each trial.

\subsection{RMS Amplitude}

Root-mean-square amplitude measures the overall signal magnitude:

\begin{equation}
RMS_i^{(k)} =
\sqrt{
\frac{1}{T}
\sum_{t}
x_i^{(k)}(t)^2
}
\end{equation}

Large RMS values may indicate noise bursts or artifacts.

\subsection{Variance}

Signal variance measures fluctuation amplitude:

\begin{equation}
Var_i^{(k)} =
\frac{1}{T}
\sum_{t}
(x_i^{(k)}(t) - \bar{x}_i^{(k)})^2
\end{equation}

where $\bar{x}_i^{(k)}$ is the mean baseline signal.

\subsection{Peak-to-Peak Amplitude}

Peak-to-peak amplitude captures extreme signal excursions:

\begin{equation}
P2P_i^{(k)} =
\max_t x_i^{(k)}(t) -
\min_t x_i^{(k)}(t)
\end{equation}

Large peak-to-peak values often indicate transient artifacts.

\subsection{Kurtosis}

Kurtosis measures deviations from Gaussian amplitude distributions:

\begin{equation}
Kurt_i^{(k)} =
\frac{E[(x - \mu)^4]}{\sigma^4}
\end{equation}

High kurtosis may indicate rare but extreme signal deviations.

\section{Composite Quality Score}

To combine multiple metrics into a single quality measure,
each statistic is normalized relative to the distribution
across trials.

Let $z_m^{(k)}$ denote the z-scored value of metric $m$ for trial $k$.

A composite quality score is then defined as

\begin{equation}
Q^{(k)} =
\sum_m w_m z_m^{(k)}
\end{equation}

where $w_m$ are weighting coefficients.

This score summarizes the likelihood that a trial contains
abnormal signal activity.

\section{Trial Classification}

Trials are categorized according to their composite quality score.

Typical classification scheme:

\begin{itemize}
\item \textbf{Green trials}: high-quality trials retained for analysis
\item \textbf{Yellow trials}: moderate artifacts, down-weighted
\item \textbf{Red trials}: severe artifacts, excluded
\end{itemize}

Trials labeled as red are excluded from ERP estimation and
time--frequency analysis.

\chapter{Event-Related Potential Analysis}

Event-related potentials (ERPs) provide a classical representation
of stimulus-locked neural responses. ERPs represent the average
neural activity across repeated trials aligned to an experimental
event.

\section{Robust ERP Estimation}

Let $x_i^{(k)}(t)$ denote the signal from channel $i$ during trial $k$.

The ERP is computed as the average across trials:

\begin{equation}
ERP_i(t) =
\frac{1}{K}
\sum_{k=1}^{K}
x_i^{(k)}(t)
\end{equation}

However, simple averaging can be sensitive to outlier trials.
Robust ERP estimation therefore incorporates trial weighting
or robust statistics.

\section{Weighted Averaging}

Let $w_k$ denote the weight assigned to trial $k$ based on
its quality score.

The weighted ERP is defined as

\begin{equation}
ERP_i(t) =
\frac{
\sum_{k} w_k x_i^{(k)}(t)
}{
\sum_{k} w_k
}
\end{equation}

where trials with lower quality receive smaller weights.

\section{ERP Component Extraction}

ERP waveforms often contain characteristic components
reflecting neural processing stages.

These components are typically characterized by:

\begin{itemize}
\item peak amplitude
\item peak latency
\item component width
\end{itemize}

Peak amplitude is defined as

\begin{equation}
A_i =
\max_{t \in T_{ROI}} ERP_i(t)
\end{equation}

where $T_{ROI}$ denotes a predefined time window of interest.

ERP component features provide summary measures of neural
responses that can be used for statistical comparisons across
conditions.




\chapter{Time--Frequency Analysis}

Neural signals recorded using SEEG exhibit complex oscillatory
dynamics that evolve over time. Classical Fourier analysis provides
information about frequency content but does not capture temporal
changes in spectral activity. Time--frequency analysis therefore
provides a framework for estimating spectral power as a function
of both time and frequency.

In this work, time--frequency representations are computed using
Morlet wavelet transforms, which provide a flexible method for
estimating oscillatory activity across a wide range of frequencies.

\section{Time--Frequency Representation}

A neural signal $x(t)$ can be represented in the time--frequency
domain as

\begin{equation}
X(t,f)
\end{equation}

where $X(t,f)$ represents the complex spectral coefficient
at time $t$ and frequency $f$.

Time--frequency analysis therefore transforms a one-dimensional
time series

\begin{equation}
x(t)
\end{equation}

into a two-dimensional representation

\begin{equation}
X(t,f)
\end{equation}

that describes how spectral power evolves over time.

\section{Morlet Wavelet Transform}

The Morlet wavelet is widely used for neural signal analysis due
to its favorable time--frequency localization properties.

The complex Morlet wavelet is defined as

\begin{equation}
\psi(t,f) =
e^{2\pi i f t}
\,
e^{-\frac{t^2}{2\sigma^2}}
\end{equation}

where

\begin{itemize}
\item $f$ is the center frequency
\item $\sigma$ controls the temporal width of the wavelet
\end{itemize}

The Gaussian envelope ensures that the wavelet is localized in time.

The relationship between the temporal width $\sigma$ and the number
of wavelet cycles $n_c$ is

\begin{equation}
\sigma = \frac{n_c}{2\pi f}
\end{equation}

Typical analyses use

\begin{equation}
n_c = 5\text{--}7
\end{equation}

cycles.

\section{Wavelet Convolution}

The time--frequency representation of a signal is obtained by
convolving the signal with the complex wavelet.

\begin{equation}
Z(t,f) =
x(t) * \psi(t,f)
\end{equation}

where $*$ denotes convolution.

The result is a complex-valued signal

\begin{equation}
Z(t,f) = A(t,f)e^{i\phi(t,f)}
\end{equation}

where

\begin{itemize}
\item $A(t,f)$ represents amplitude
\item $\phi(t,f)$ represents phase
\end{itemize}

\section{Spectral Power Estimation}

Spectral power is obtained from the magnitude of the complex
wavelet coefficient:

\begin{equation}
P(t,f) =
|Z(t,f)|^2
\end{equation}

This quantity represents the instantaneous power of oscillatory
activity at frequency $f$ and time $t$.

\section{Baseline Normalization}

To interpret event-related spectral changes, spectral power is
normalized relative to a prestimulus baseline period.

Let $P_{base}(f)$ denote the average power during the baseline
interval

\begin{equation}
P_{base}(f) =
\frac{1}{T_b}
\sum_{t \in \text{baseline}} P(t,f)
\end{equation}

Baseline-normalized power can then be expressed in decibel units:

\begin{equation}
P_{dB}(t,f) =
10\log_{10}
\left(
\frac{P(t,f)}{P_{base}(f)}
\right)
\end{equation}

This normalization highlights relative increases or decreases
in oscillatory power following experimental events.

\section{Evoked and Induced Oscillations}

Time--frequency responses can be separated into evoked and induced
components.

Evoked activity is phase-locked to the stimulus and is reflected
in the event-related potential (ERP). Induced activity is not
phase-locked and therefore varies across trials.

To isolate induced oscillations, the ERP is first subtracted from
each trial:

\begin{equation}
x_{induced}^{(k)}(t) =
x^{(k)}(t) - ERP(t)
\end{equation}

where $k$ indexes trials.

Wavelet analysis is then applied to the residual signal
$x_{induced}^{(k)}(t)$ to estimate induced spectral power.

\section{Inter-Trial Phase Coherence}

Phase-based synchronization across trials can be quantified using
inter-trial phase coherence (ITPC).

Let $\phi_k(t,f)$ denote the phase of the wavelet coefficient for
trial $k$.

ITPC is defined as

\begin{equation}
ITPC(t,f) =
\left|
\frac{1}{N}
\sum_{k=1}^{N}
e^{i\phi_k(t,f)}
\right|
\end{equation}

where $N$ is the number of trials.

ITPC values range between

\begin{itemize}
\item $0$ (random phase across trials)
\item $1$ (perfect phase locking)
\end{itemize}

ITPC therefore measures the consistency of phase alignment
across repeated trials.

\section{Frequency Bands of Interest}

Neural oscillations are often summarized within canonical frequency
bands:

\begin{center}
\begin{tabular}{ll}
Theta & 4--8 Hz \\
Alpha & 8--12 Hz \\
Beta & 13--30 Hz \\
Gamma & 30--70 Hz \\
High-frequency activity & 70--150 Hz
\end{tabular}
\end{center}

Band-limited power can be obtained by averaging spectral power
across frequencies within each band.

\section{Output of Time--Frequency Analysis}

The time--frequency analysis produces the following outputs:

\begin{itemize}
\item Spectral power $P(t,f)$
\item Baseline-normalized power $P_{dB}(t,f)$
\item Induced spectral power
\item Inter-trial phase coherence (ITPC)
\item Band-limited power timecourses
\end{itemize}

These representations form the basis for subsequent statistical
analysis and neural population dynamics modeling.


\chapter{Oscillatory Neural Dynamics}

Neural activity recorded using SEEG exhibits oscillatory dynamics across
multiple frequency bands. Oscillations reflect coordinated activity
of neural populations and are commonly modulated during sensory,
cognitive, and motor processes.

Oscillatory activity can be broadly divided into two components:

\begin{itemize}
\item Evoked oscillations, which are phase-locked to the stimulus
\item Induced oscillations, which are not phase-locked and vary across trials
\end{itemize}

This chapter focuses on induced oscillatory activity and broadband
high-frequency responses.

\section{Induced Oscillations}

Induced oscillations represent changes in spectral power that are not
phase-locked to the stimulus. These responses vary across trials and
are therefore not visible in the averaged ERP waveform.

Let $x^{(k)}(t)$ denote the signal from trial $k$ and
$ERP(t)$ the averaged event-related potential.

Induced signals are obtained by removing the ERP from each trial:

\begin{equation}
x^{(k)}_{ind}(t) =
x^{(k)}(t) - ERP(t)
\end{equation}

Time--frequency analysis is then applied to the residual signal to
estimate induced spectral power.

\section{High-Frequency Activity}

High-frequency activity (HFA), typically in the range
70--150 Hz, is widely used as a proxy for local neuronal population
firing.

Let $P(t,f)$ denote spectral power obtained from the wavelet transform.
High-frequency activity is computed by averaging power across the
high-frequency band:

\begin{equation}
HFA(t) =
\frac{1}{|F|}
\sum_{f \in F} P(t,f)
\end{equation}

where $F$ denotes the set of frequencies within the high-frequency band.

High-frequency activity provides a sensitive measure of local neural
activation.

\section{Band-Limited Power}

Spectral power is often summarized within canonical frequency bands.
Band-limited power is obtained by averaging spectral power across
frequencies within a given band.

\begin{equation}
B_{band}(t) =
\frac{1}{|F_{band}|}
\sum_{f \in F_{band}} P(t,f)
\end{equation}

Typical frequency bands include:

\begin{center}
\begin{tabular}{ll}
Theta & 4--8 Hz \\
Alpha & 8--12 Hz \\
Beta & 13--30 Hz \\
Gamma & 30--70 Hz \\
High-frequency activity & 70--150 Hz
\end{tabular}
\end{center}

Band-limited power timecourses provide compact representations of
oscillatory neural responses.

\chapter{Phase-Based Neural Synchronization}

In addition to spectral power, neural oscillations contain phase
information that reflects synchronization across trials or brain
regions.

Phase-based metrics capture stimulus-locked phase alignment of
oscillatory activity.

\section{Inter-Trial Phase Coherence}

Inter-trial phase coherence (ITPC) measures the consistency of phase
across trials.

Let $\phi_k(t,f)$ denote the phase of the complex wavelet coefficient
for trial $k$.

ITPC is defined as

\begin{equation}
ITPC(t,f) =
\left|
\frac{1}{N}
\sum_{k=1}^{N}
e^{i\phi_k(t,f)}
\right|
\end{equation}

where $N$ is the number of trials.

ITPC ranges between 0 and 1, where

\begin{itemize}
\item $0$ indicates random phase across trials
\item $1$ indicates perfect phase locking
\end{itemize}

\section{Phase Locking Value}

Phase locking value (PLV) measures synchronization between two
recording channels.

Let $\phi_i(t,f)$ and $\phi_j(t,f)$ denote the phases of channels
$i$ and $j$.

The PLV is defined as

\begin{equation}
PLV_{ij}(t,f) =
\left|
\frac{1}{N}
\sum_{k=1}^{N}
e^{i(\phi_{i,k}(t,f)-\phi_{j,k}(t,f))}
\right|
\end{equation}

PLV captures phase coupling between spatially distributed neural
populations.

\chapter{Statistical Inference}

Statistical inference is used to determine whether observed neural
responses differ significantly from baseline activity or across
experimental conditions.

Because neural data often violate parametric assumptions,
nonparametric permutation tests are widely used.

\section{Permutation Testing}

Permutation testing constructs a null distribution by randomly
reassigning condition labels across trials.

Let $T$ denote a test statistic computed from the observed data.
Permutation testing proceeds as follows:

\begin{enumerate}
\item Randomly shuffle condition labels
\item Recompute the statistic $T$
\item Repeat for many permutations
\end{enumerate}

The empirical p-value is computed as

\begin{equation}
p =
\frac{\#(T_{perm} \ge T_{obs}) + 1}{N_{perm} + 1}
\end{equation}

where $T_{perm}$ are statistics obtained from permuted data.

\section{ROI-Based Statistics}

Region-of-interest (ROI) statistics summarize neural responses within
predefined time or frequency windows.

For example, band-limited power may be averaged over a time interval

\begin{equation}
S =
\frac{1}{|T|}
\sum_{t \in T} B_{band}(t)
\end{equation}

Permutation testing is then applied to compare this statistic across
conditions.

\section{Cluster-Based Time--Frequency Statistics}

Cluster-based permutation tests control the multiple-comparison
problem inherent in time--frequency maps.

Adjacent time--frequency points exceeding a statistical threshold
are grouped into clusters.

The cluster statistic is defined as

\begin{equation}
C =
\sum_{(t,f) \in cluster} T(t,f)
\end{equation}

where $T(t,f)$ denotes the local test statistic.

Cluster significance is evaluated using permutation-based null
distributions.

\section{Multiple Comparison Correction}

Because SEEG analyses often involve many channels and timepoints,
multiple comparison correction is necessary.

False discovery rate (FDR) correction controls the expected
proportion of false positives among significant tests.

\section{Response Onset Detection}

Response onset detection identifies the earliest timepoint at which
neural activity deviates significantly from baseline.

Permutation-based max-statistic methods are used to control
family-wise error across timepoints.

\chapter{Multichannel Neural Population Representation}

Single-channel analysis captures local neural responses but does not
fully characterize distributed neural dynamics.

Multichannel SEEG recordings are therefore represented as neural
population states.

\section{Feature Tensor Construction}

Neural features extracted from each channel are combined into a
multidimensional tensor representation.

Let

\begin{equation}
X(k,i,t)
\end{equation}

denote the neural feature from trial $k$, channel $i$, and time $t$.

This representation captures the spatiotemporal structure of neural
population activity.

\section{Trial × Channel × Time Representations}

The resulting tensor

\begin{equation}
X \in \mathbb{R}^{K \times N \times T}
\end{equation}

contains

\begin{itemize}
\item $K$ trials
\item $N$ channels
\item $T$ time samples
\end{itemize}

For population-level analyses, this tensor is reshaped into a
matrix representation

\begin{equation}
X(t) \in \mathbb{R}^{N}
\end{equation}

representing the neural population state at time $t$.

\chapter{Dimensionality Reduction}

Multichannel neural population activity often evolves within a
low-dimensional subspace of the high-dimensional recording space.

Dimensionality reduction methods identify latent variables that
capture coordinated neural dynamics.

\section{Principal Component Analysis}

Principal component analysis (PCA) identifies orthogonal directions
that maximize variance in the neural data.

Let $X$ denote the data matrix with dimensions $N \times T$.

PCA computes eigenvectors of the covariance matrix

\begin{equation}
\Sigma =
\frac{1}{T} XX^T
\end{equation}

The principal components correspond to eigenvectors of $\Sigma$.

Latent neural states are obtained by projecting the data onto the
leading components:

\begin{equation}
z(t) = W^T x(t)
\end{equation}

where $W$ contains the principal component vectors.

\section{Factor Analysis}

Factor analysis models neural activity as a combination of latent
factors and independent noise.

\begin{equation}
x(t) = W z(t) + \epsilon
\end{equation}

where

\begin{itemize}
\item $z(t)$ represents latent factors
\item $W$ is a loading matrix
\item $\epsilon$ is Gaussian noise
\end{itemize}

Factor analysis separates shared neural variability from
channel-specific noise.

\section{Gaussian Process Factor Analysis}

Gaussian process factor analysis (GPFA) extends factor analysis
by modeling latent trajectories as smooth temporal processes.

Each latent dimension follows a Gaussian process prior:

\begin{equation}
z_d(t) \sim GP(0, k(t,t'))
\end{equation}

where $k(t,t')$ is a covariance kernel that enforces temporal
smoothness.

GPFA is particularly useful for recovering smooth neural trajectories
from noisy population recordings.


\chapter{Latent Neural Trajectories}

Multichannel SEEG recordings capture high-dimensional neural activity
across distributed brain regions. However, neural population activity
often evolves within a low-dimensional subspace. Dimensionality
reduction methods such as PCA or factor analysis can therefore
identify latent neural variables that summarize coordinated neural
dynamics.

Let $x(t) \in \mathbb{R}^{N}$ denote the multichannel neural activity
vector at time $t$, where $N$ is the number of recording channels.

Dimensionality reduction projects the neural activity into a
low-dimensional latent state space:

\begin{equation}
z(t) = W^{T}x(t)
\end{equation}

where

\begin{itemize}
\item $z(t) \in \mathbb{R}^{D}$ is the latent neural state
\item $W$ is the projection matrix
\item $D \ll N$ is the dimensionality of the latent space
\end{itemize}

The time evolution of the latent state $z(t)$ forms a neural trajectory
within the low-dimensional state space.

\section{Single-Trial Trajectories}

For each trial $k$, the multichannel neural activity can be represented
as a time-varying latent state vector:

\begin{equation}
z^{(k)}(t) = W^{T} x^{(k)}(t)
\end{equation}

where

\begin{itemize}
\item $x^{(k)}(t)$ is the neural population vector during trial $k$
\item $z^{(k)}(t)$ is the corresponding latent neural state
\end{itemize}

The sequence of latent states across time forms a trajectory:

\begin{equation}
\mathbf{z}^{(k)} =
\left\{
z^{(k)}(t_1),
z^{(k)}(t_2),
\dots,
z^{(k)}(t_T)
\right\}
\end{equation}

Each trial therefore corresponds to a trajectory through the latent
state space.

These trajectories describe how neural population activity evolves
over time in response to experimental events.

Because neural recordings are noisy, trajectories are often smoothed
using temporal filtering. For example, Gaussian smoothing can be
applied to each latent dimension:

\begin{equation}
\tilde{z}^{(k)}(t) =
\sum_{\tau} z^{(k)}(\tau)
\, g(t-\tau)
\end{equation}

where $g(t)$ denotes a Gaussian smoothing kernel.

Single-trial trajectories provide insight into trial-to-trial
variability of neural population dynamics.

\section{Condition-Averaged Trajectories}

To characterize population responses associated with different
experimental conditions, trajectories are averaged across trials
belonging to the same condition.

Let $\mathcal{C}$ denote the set of trials belonging to a particular
condition.

The condition-averaged trajectory is defined as

\begin{equation}
\bar{z}_{\mathcal{C}}(t) =
\frac{1}{|\mathcal{C}|}
\sum_{k \in \mathcal{C}}
z^{(k)}(t)
\end{equation}

This trajectory represents the average neural population response
for that condition.

Condition-averaged trajectories can be visualized in two- or
three-dimensional latent space to reveal differences between
experimental conditions.

In addition, the distance between condition trajectories can be
computed as

\begin{equation}
d(t) =
\left\|
\bar{z}_{A}(t) - \bar{z}_{B}(t)
\right\|
\end{equation}

where $A$ and $B$ denote two experimental conditions.

Trajectory separation over time provides a quantitative measure of
how neural population dynamics differ between conditions.

The latent neural trajectories obtained in this chapter serve as the
basis for subsequent analyses of neural population geometry and
dynamical systems structure.


\chapter{Trajectory Geometry}

Latent neural trajectories describe how population activity evolves
through a low-dimensional state space. The geometric properties of
these trajectories provide insight into the structure of neural
population dynamics.

Let

\begin{equation}
z(t) \in \mathbb{R}^{D}
\end{equation}

denote the latent neural state at time $t$, where $D$ is the number
of latent dimensions.

The sequence

\begin{equation}
\mathbf{z} = \{z(t_1), z(t_2), \dots, z(t_T)\}
\end{equation}

defines a trajectory through latent state space.

Geometric properties of this trajectory can be quantified using
velocity, curvature, dispersion, and path length.

\section{Trajectory Velocity}

Trajectory velocity describes the rate at which the neural population
state changes over time.

The instantaneous velocity is defined as the temporal derivative of
the latent state:

\begin{equation}
v(t) =
\frac{d z(t)}{dt}
\end{equation}

In discrete time, velocity is approximated using finite differences:

\begin{equation}
v(t) =
z(t+\Delta t) - z(t)
\end{equation}

The magnitude of the velocity vector provides a scalar measure of
the speed of neural state evolution:

\begin{equation}
||v(t)|| =
\sqrt{\sum_{d=1}^{D} v_d(t)^2}
\end{equation}

High velocity indicates rapid transitions in neural population
activity, whereas low velocity corresponds to relatively stable
neural states.

\section{Trajectory Curvature}

Trajectory curvature quantifies how sharply the neural trajectory
changes direction in latent state space.

Let

\begin{equation}
v(t) = \frac{dz}{dt}
\end{equation}

denote the velocity vector, and

\begin{equation}
a(t) = \frac{d^2 z}{dt^2}
\end{equation}

denote the acceleration vector.

Curvature is defined as

\begin{equation}
\kappa(t) =
\frac{||v(t) \times a(t)||}{||v(t)||^3}
\end{equation}

In practice, curvature can be approximated using discrete derivatives
of the latent trajectory.

High curvature indicates rapid changes in trajectory direction,
suggesting transitions between neural population states.

\section{Trajectory Dispersion}

Trajectory dispersion measures the variability of neural trajectories
across trials.

Let $z^{(k)}(t)$ denote the latent trajectory for trial $k$, and
$\bar{z}(t)$ the condition-averaged trajectory.

Dispersion is defined as the average distance between single-trial
trajectories and the mean trajectory:

\begin{equation}
D(t) =
\frac{1}{K}
\sum_{k=1}^{K}
||z^{(k)}(t) - \bar{z}(t)||
\end{equation}

where $K$ is the number of trials.

High dispersion indicates greater trial-to-trial variability in neural
population dynamics.

\section{Path Length}

Trajectory path length measures the total distance traveled by the
neural trajectory in latent state space.

Let

\begin{equation}
v(t) = z(t+\Delta t) - z(t)
\end{equation}

denote the incremental displacement vector.

The total path length is defined as

\begin{equation}
L =
\sum_{t=1}^{T-1}
||z(t+1) - z(t)||
\end{equation}

Path length provides a measure of the overall magnitude of neural
population state changes during the response period.

Longer path lengths indicate greater neural population dynamics
following experimental events.

The geometric metrics defined in this chapter provide quantitative
descriptions of neural population trajectories and form the basis for
subsequent analyses of condition separation and dynamical systems
structure.


\chapter{Condition Separation in Latent Space}

Latent neural trajectories provide a compact representation of
multichannel neural population dynamics. When neural responses differ
between experimental conditions, the corresponding trajectories may
diverge within the latent state space.

Quantifying this divergence allows identification of neural population
dynamics that distinguish experimental conditions.

Let

\begin{equation}
z^{(k)}(t) \in \mathbb{R}^{D}
\end{equation}

denote the latent neural state at time $t$ for trial $k$.

Trials are grouped according to experimental condition
(e.g., condition $A$ and condition $B$).

\section{Condition Centroids}

For each condition, the average latent trajectory is computed across
all trials belonging to that condition.

Let $\mathcal{C}$ denote the set of trials for a given condition.

The condition centroid trajectory is defined as

\begin{equation}
\bar{z}_{\mathcal{C}}(t) =
\frac{1}{|\mathcal{C}|}
\sum_{k \in \mathcal{C}} z^{(k)}(t)
\end{equation}

This trajectory represents the average neural population response
associated with that experimental condition.

Condition centroids provide a simplified representation of neural
population dynamics and are commonly used to visualize condition
differences in latent state space.

\section{Inter-Condition Distance}

To quantify differences between conditions, the distance between
condition centroid trajectories can be computed.

The simplest measure is Euclidean distance:

\begin{equation}
d(t) =
\left\|
\bar{z}_{A}(t) -
\bar{z}_{B}(t)
\right\|
\end{equation}

where $\bar{z}_{A}(t)$ and $\bar{z}_{B}(t)$ denote centroid trajectories
for conditions $A$ and $B$.

Alternatively, Mahalanobis distance can be used to account for
covariance structure of the neural state space:

\begin{equation}
d_M(t) =
\sqrt{
(\bar{z}_A(t)-\bar{z}_B(t))^T
\Sigma^{-1}
(\bar{z}_A(t)-\bar{z}_B(t))
}
\end{equation}

where $\Sigma$ is the covariance matrix of latent states.

Mahalanobis distance provides a scale-invariant measure of
condition separation.

\section{Time-Resolved Trajectory Separation}

Condition separation is often evaluated as a function of time in
order to identify when neural trajectories begin to diverge.

For each timepoint $t$, a separation statistic is computed using
the distances between single-trial trajectories and condition
centroids.

For example, a separation index can be defined as

\begin{equation}
S(t) =
\frac{
\left\|
\bar{z}_{A}(t) - \bar{z}_{B}(t)
\right\|
}{
\sigma_A(t) + \sigma_B(t)
}
\end{equation}

where

\begin{equation}
\sigma_A(t) =
\frac{1}{|A|}
\sum_{k \in A}
\left\|
z^{(k)}(t) - \bar{z}_A(t)
\right\|
\end{equation}

measures within-condition dispersion.

Large values of $S(t)$ indicate strong separation between
condition trajectories.

Statistical significance of trajectory separation can be assessed
using permutation testing. Condition labels are randomly shuffled
across trials, and the separation statistic is recomputed to
generate a null distribution.

The observed statistic is then compared to this null distribution
to determine whether neural trajectories diverge significantly
between conditions.

Timepoints at which significant separation occurs indicate when
neural population dynamics begin to encode experimental condition
information.

\chapter{Rotational Neural Dynamics}

Neural population activity often evolves according to structured
dynamical patterns. In many neural systems, population responses
exhibit rotational structure in low-dimensional latent state space.
Such rotational dynamics have been observed in motor cortex,
sensory systems, and other neural populations.

Rotational dynamics reflect coordinated activity patterns that
evolve according to underlying neural circuits.

This chapter introduces methods for identifying rotational
structure in latent neural trajectories using jPCA.

\section{Neural Population Dynamics}

Let

\begin{equation}
z(t) \in \mathbb{R}^{D}
\end{equation}

denote the latent neural state vector obtained from dimensionality
reduction of multichannel SEEG activity.

The temporal evolution of neural population activity can be
approximated using a linear dynamical system:

\begin{equation}
\frac{dz(t)}{dt} = M z(t)
\end{equation}

where

\begin{itemize}
\item $z(t)$ is the latent neural state
\item $M$ is a dynamics matrix
\end{itemize}

The matrix $M$ determines how neural states evolve through time.

\section{Rotational Dynamics}

If the dynamics matrix $M$ is skew-symmetric

\begin{equation}
M = -M^T
\end{equation}

then the resulting dynamics correspond to pure rotations in state
space.

In this case,

\begin{equation}
\frac{dz(t)}{dt} = M z(t)
\end{equation}

generates circular trajectories in the latent space.

Such rotational dynamics indicate that neural activity evolves
through structured population states rather than simply increasing
or decreasing in amplitude.

\section{jPCA Analysis}

jPCA (Jordan canonical projection analysis) is used to identify
rotational structure in neural population activity.

The method proceeds by estimating the temporal derivative of the
latent trajectory:

\begin{equation}
\dot{z}(t) =
\frac{dz(t)}{dt}
\end{equation}

A linear model is then fitted:

\begin{equation}
\dot{z}(t) = M z(t)
\end{equation}

The matrix $M$ is constrained to be skew-symmetric:

\begin{equation}
M = -M^T
\end{equation}

This constraint ensures that the resulting dynamics correspond
to rotational flow.

The skew-symmetric matrix can be parameterized as

\begin{equation}
M =
\begin{bmatrix}
0 & -a \\
a & 0
\end{bmatrix}
\end{equation}

for a two-dimensional rotational plane.

\section{Extraction of Rotational Planes}

Eigenvectors of the dynamics matrix define rotational planes in
the latent state space.

The dominant rotational components correspond to eigenvectors
associated with the largest imaginary eigenvalues.

These components define a two-dimensional plane in which neural
trajectories exhibit the strongest rotational structure.

Let $W_{rot}$ denote the projection matrix for the rotational
subspace.

Rotational trajectories are then obtained as

\begin{equation}
z_{rot}(t) = W_{rot}^{T} z(t)
\end{equation}

\section{Visualization of Rotational Dynamics}

Rotational neural dynamics can be visualized by plotting neural
trajectories in the rotational plane.

For each experimental condition, the condition-averaged trajectory

\begin{equation}
\bar{z}(t)
\end{equation}

is projected into the rotational subspace:

\begin{equation}
z_{rot}(t) = W_{rot}^{T} \bar{z}(t)
\end{equation}

The resulting two-dimensional trajectory illustrates the temporal
evolution of neural population activity.

Differences between experimental conditions may appear as distinct
rotational trajectories or phase shifts within the rotational
plane.

Identification of rotational dynamics provides evidence that
neural population responses are governed by low-dimensional
dynamical systems rather than independent channel responses.


\chapter{Tangent Space Dynamics}

Neural population trajectories describe the evolution of latent neural
states over time. While trajectory geometry captures global properties
such as velocity and curvature, the local dynamics of neural activity
can be characterized by tangent vectors along the trajectory.

Tangent vectors represent the instantaneous direction of neural
population state evolution and therefore describe the local flow of
neural dynamics in latent state space.

Analyzing tangent vectors allows comparison of neural dynamics across
experimental conditions and provides insight into the structure of the
underlying neural dynamical system.

\section{Tangent Vector Estimation}

Let

\begin{equation}
z(t) \in \mathbb{R}^{D}
\end{equation}

denote the latent neural state vector at time $t$.

The tangent vector of the trajectory is defined as the temporal
derivative of the latent state:

\begin{equation}
v(t) = \frac{dz(t)}{dt}
\end{equation}

In discrete time, tangent vectors can be approximated using finite
differences:

\begin{equation}
v(t) =
\frac{z(t+\Delta t) - z(t-\Delta t)}{2\Delta t}
\end{equation}

where $\Delta t$ is the sampling interval.

The resulting vector

\begin{equation}
v(t) \in \mathbb{R}^{D}
\end{equation}

represents the local direction of neural population flow in the latent
state space.

For each trial $k$, tangent vectors can be estimated as

\begin{equation}
v^{(k)}(t) =
\frac{z^{(k)}(t+\Delta t) - z^{(k)}(t-\Delta t)}{2\Delta t}
\end{equation}

These vectors characterize the instantaneous dynamics of neural
population activity.

To reduce noise in derivative estimates, tangent vectors are often
computed after temporal smoothing of the latent trajectories.

\section{Principal Angle Analysis}

To compare neural dynamics across experimental conditions, tangent
vectors can be analyzed as subspaces in latent state space.

Let

\begin{equation}
V_A(t) =
[v_A^{(1)}(t), v_A^{(2)}(t), \dots]
\end{equation}

denote the set of tangent vectors for condition $A$, and

\begin{equation}
V_B(t)
\end{equation}

the corresponding vectors for condition $B$.

These vectors span tangent subspaces that describe the local flow of
neural dynamics for each condition.

Similarity between tangent subspaces can be quantified using
principal angles.

Let $U_A$ and $U_B$ denote orthonormal bases for the tangent
subspaces of conditions $A$ and $B$.

Principal angles $\theta_i$ are defined as

\begin{equation}
\cos(\theta_i) =
\sigma_i(U_A^T U_B)
\end{equation}

where $\sigma_i$ denotes the singular values obtained from
singular value decomposition.

Small principal angles indicate similar neural dynamics across
conditions, whereas large angles indicate divergent population
flows.

Principal angle analysis therefore provides a geometric measure
of differences in neural population dynamics beyond simple
trajectory separation.

Time-resolved principal angle analysis can reveal when neural
population dynamics begin to diverge between experimental
conditions.

\chapter{Geometry of Neural Manifolds}

Neural population activity often evolves within a low-dimensional
subspace embedded in the high-dimensional neural recording space.
This low-dimensional structure is commonly referred to as a
\textit{neural manifold}. The neural manifold captures coordinated
patterns of activity across multiple channels and reflects the
collective dynamics of neural populations.

Let

\begin{equation}
x(t) \in \mathbb{R}^{N}
\end{equation}

denote the multichannel neural activity vector, where $N$ is the
number of recording channels. After dimensionality reduction, the
latent neural state is represented as

\begin{equation}
z(t) \in \mathbb{R}^{D}, \quad D \ll N
\end{equation}

where $D$ denotes the dimensionality of the latent state space.

The neural manifold is defined as the subset of latent space visited
by neural trajectories during task performance.

\section{State-Space Occupancy}

State-space occupancy describes the region of latent space explored
by neural population activity.

Let

\begin{equation}
\mathcal{Z} = \{ z(t_1), z(t_2), \dots, z(t_T) \}
\end{equation}

denote the set of latent neural states observed across time.

A simple measure of state-space occupancy is the covariance matrix
of latent states:

\begin{equation}
\Sigma =
\frac{1}{T}
\sum_{t=1}^{T}
(z(t) - \bar{z})
(z(t) - \bar{z})^T
\end{equation}

where $\bar{z}$ denotes the mean latent state.

Eigenvalues of $\Sigma$ characterize the spread of neural activity
along different latent dimensions.

Large eigenvalues indicate directions of strong neural population
variability, whereas small eigenvalues correspond to constrained
dimensions of the neural manifold.

\section{Trajectory Recurrence}

Trajectory recurrence measures the tendency of neural population
states to revisit previously occupied regions of latent space.

For two timepoints $t_i$ and $t_j$, the distance between neural
states is defined as

\begin{equation}
d_{ij} =
\| z(t_i) - z(t_j) \|
\end{equation}

A recurrence event occurs when

\begin{equation}
d_{ij} < \epsilon
\end{equation}

where $\epsilon$ is a small threshold.

Recurrence matrices can be constructed as

\begin{equation}
R_{ij} =
\begin{cases}
1 & \text{if } d_{ij} < \epsilon \\
0 & \text{otherwise}
\end{cases}
\end{equation}

This matrix describes when neural population states return to
previous regions of the manifold.

Patterns in the recurrence matrix can reveal periodic dynamics
or attractor-like structure in neural population activity.

\section{Subspace Overlap}

Neural manifolds associated with different experimental conditions
can be compared by measuring the overlap between their respective
subspaces.

Let

\begin{equation}
U_A
\end{equation}

and

\begin{equation}
U_B
\end{equation}

denote orthonormal basis matrices spanning the neural manifolds
for conditions $A$ and $B$.

The overlap between these subspaces can be quantified using
principal angles.

Let $\theta_i$ denote the principal angles defined by

\begin{equation}
\cos(\theta_i) =
\sigma_i(U_A^T U_B)
\end{equation}

where $\sigma_i$ are singular values obtained from singular value
decomposition.

Small principal angles indicate strong overlap between neural
manifolds, whereas large angles indicate that neural population
activity occupies different subspaces under the two conditions.

Subspace overlap therefore provides a quantitative measure of
differences in the global structure of neural population activity.

\chapter{Neural Decoding}

Latent neural states extracted from multichannel SEEG recordings
provide a compact representation of neural population dynamics.
If neural activity encodes experimental variables, these latent
states should contain information that allows prediction of
experimental conditions or behavioral events.

Neural decoding methods quantify how well experimental variables
can be predicted from neural population activity.

\section{Latent-State Representation}

Let

\begin{equation}
z^{(k)}(t) \in \mathbb{R}^{D}
\end{equation}

denote the latent neural state vector for trial $k$ at time $t$,
obtained from dimensionality reduction.

Each trial is associated with a condition label

\begin{equation}
y^{(k)} \in \{1,2,\dots,C\}
\end{equation}

where $C$ denotes the number of experimental conditions.

The decoding task consists of predicting the condition label
$y^{(k)}$ from the neural state $z^{(k)}(t)$.

\section{Linear Classifier}

A common approach is to use a linear classifier that maps latent
neural states to predicted condition labels.

For example, a linear discriminant classifier computes

\begin{equation}
\hat{y} =
\arg\max_c
\left(
w_c^T z + b_c
\right)
\end{equation}

where

\begin{itemize}
\item $w_c$ is the classifier weight vector
\item $b_c$ is the bias term
\end{itemize}

Classifier parameters are estimated using training data.

\section{Cross-Validation}

To avoid overfitting, decoding performance is evaluated using
cross-validation.

Trials are partitioned into training and testing sets.
The classifier is trained on the training set and evaluated
on held-out test trials.

Typical cross-validation schemes include:

\begin{itemize}
\item k-fold cross-validation
\item leave-one-trial-out
\item leave-one-condition-out
\end{itemize}

Decoding accuracy is computed as

\begin{equation}
Accuracy =
\frac{\text{Number of correct predictions}}
{\text{Total number of test trials}}
\end{equation}

\section{Time-Resolved Decoding}

To determine when neural population activity begins to encode
experimental variables, decoding can be performed separately
for each timepoint.

For each time $t$, the classifier is trained using latent states
$z^{(k)}(t)$ and corresponding labels $y^{(k)}$.

This produces a decoding accuracy timecourse

\begin{equation}
A(t)
\end{equation}

which reflects the temporal evolution of decodable information
in the neural population.

\section{Permutation Testing}

To assess statistical significance of decoding performance,
permutation testing is used.

Condition labels are randomly shuffled across trials, and
decoding accuracy is recomputed.

This procedure generates a null distribution

\begin{equation}
A_{perm}(t)
\end{equation}

representing decoding accuracy expected by chance.

The observed decoding accuracy is compared to the null
distribution to compute a p-value.

\section{Decoding Onset Detection}

An important question is when neural activity begins to encode
experimental information.

Decoding onset is defined as the earliest timepoint at which
decoding accuracy exceeds the permutation threshold.

Formally,

\begin{equation}
t_{onset} =
\min
\left\{
t \mid
A(t) > A_{threshold}
\right\}
\end{equation}

where $A_{threshold}$ is determined from the permutation
distribution.

Decoding onset therefore identifies the time at which neural
population activity begins to represent experimental conditions.

\section{Relationship to Neural Trajectories}

Decoding results can be related to the geometry of neural
trajectories.

When trajectories corresponding to different conditions diverge
in latent state space, decoding accuracy increases.

Thus, neural decoding provides a complementary perspective on
neural population dynamics, linking latent trajectory geometry
to information encoding in neural populations.


\chapter{Uncertainty Quantification}

Neural trajectories estimated from multichannel SEEG recordings
reflect both true neural population dynamics and variability arising
from finite sampling and measurement noise. Quantifying uncertainty
in trajectory estimates is therefore essential for reliable
interpretation of neural population analyses.

Bootstrap resampling methods provide a nonparametric approach for
estimating confidence intervals for neural trajectories and derived
trajectory statistics.

\section{Bootstrap Trajectories}

Bootstrap methods estimate the variability of neural trajectories by
resampling trials with replacement.

Let

\begin{equation}
\{ z^{(1)}(t), z^{(2)}(t), \dots, z^{(K)}(t) \}
\end{equation}

denote the set of single-trial latent trajectories for a given
experimental condition.

A bootstrap sample is generated by randomly sampling $K$ trials
with replacement from the original set of trials. The condition-
averaged trajectory is then recomputed for the bootstrap sample:

\begin{equation}
\bar{z}^{*}(t) =
\frac{1}{K}
\sum_{k=1}^{K}
z^{*(k)}(t)
\end{equation}

where $z^{*(k)}(t)$ denotes trajectories selected in the bootstrap
sample.

Repeating this procedure for $B$ bootstrap iterations produces
a set of bootstrap trajectories:

\begin{equation}
\{ \bar{z}^{*(1)}(t), \bar{z}^{*(2)}(t), \dots, \bar{z}^{*(B)}(t) \}
\end{equation}

This ensemble of trajectories provides an empirical estimate of
trajectory variability.

Bootstrap samples can also be used to estimate confidence intervals
for trajectory-derived quantities such as

\begin{itemize}
\item trajectory velocity
\item trajectory curvature
\item path length
\item condition separation distance
\end{itemize}

Confidence intervals are obtained by computing the desired statistic
for each bootstrap sample and extracting percentile bounds from the
resulting distribution.

\section{Confidence Tubes}

Confidence tubes provide a geometric representation of uncertainty
around neural trajectories in latent state space.

Let

\begin{equation}
\bar{z}(t)
\end{equation}

denote the condition-averaged trajectory and

\begin{equation}
\bar{z}^{*(b)}(t)
\end{equation}

the bootstrap trajectory from bootstrap iteration $b$.

The covariance of bootstrap trajectories at time $t$ is defined as

\begin{equation}
\Sigma(t) =
\frac{1}{B}
\sum_{b=1}^{B}
(\bar{z}^{*(b)}(t) - \bar{z}(t))
(\bar{z}^{*(b)}(t) - \bar{z}(t))^T
\end{equation}

This covariance matrix describes the uncertainty of the neural state
estimate at time $t$.

A confidence region around the trajectory can then be defined as

\begin{equation}
(z - \bar{z}(t))^T \Sigma(t)^{-1} (z - \bar{z}(t)) \le \chi^2_{\alpha}
\end{equation}

where $\chi^2_{\alpha}$ denotes the chi-squared threshold corresponding
to the desired confidence level.

The resulting region forms a confidence tube surrounding the mean
trajectory.

Confidence tubes allow visualization of uncertainty in neural
population dynamics and provide a statistical framework for
comparing trajectories across experimental conditions.

When confidence tubes for two conditions do not overlap, this
provides evidence that neural population trajectories differ
significantly between conditions.


\chapter{Visualization of Neural Dynamics}

Visualization plays a critical role in interpreting neural population
dynamics derived from SEEG recordings. Effective visualization methods
allow researchers to explore neural responses at multiple levels of
analysis, including single-channel responses, oscillatory activity,
and population-level neural trajectories.

This chapter describes visualization techniques used to represent
event-related responses, time--frequency dynamics, and neural
population trajectories in latent state space.

\section{ERP Visualization}

Event-related potentials (ERPs) provide a classical representation of
stimulus-locked neural responses.

For each channel $i$, the ERP waveform is defined as

\begin{equation}
ERP_i(t) =
\frac{1}{K}
\sum_{k=1}^{K}
x_i^{(k)}(t)
\end{equation}

where $x_i^{(k)}(t)$ denotes the signal from trial $k$.

ERP visualization typically includes the following elements:

\begin{itemize}
\item Mean ERP waveform across trials
\item Shaded confidence intervals representing trial variability
\item Event markers indicating stimulus onset
\item Comparison of ERPs across experimental conditions
\end{itemize}

Confidence intervals are commonly obtained using bootstrap
resampling of trials.

ERP plots allow identification of characteristic response components
such as early sensory responses and late cognitive components.

\section{Time--Frequency Maps}

Time--frequency maps visualize oscillatory activity across time
and frequency.

Let $P(t,f)$ denote the spectral power obtained from wavelet
transforms.

Time--frequency maps display the quantity

\begin{equation}
P_{dB}(t,f) =
10\log_{10}
\left(
\frac{P(t,f)}{P_{base}(f)}
\right)
\end{equation}

which represents baseline-normalized spectral power.

Time--frequency maps are typically visualized as heatmaps with

\begin{itemize}
\item time on the horizontal axis
\item frequency on the vertical axis
\item color representing spectral power
\end{itemize}

Statistically significant regions obtained from cluster-based
permutation tests can be overlaid on these maps to highlight
reliable oscillatory responses.

\section{Latent Trajectory Visualization}

Latent neural trajectories provide a low-dimensional representation
of multichannel neural population dynamics.

For each trial $k$, the latent trajectory is defined as

\begin{equation}
z^{(k)}(t)
\end{equation}

where $z(t)$ represents the latent neural state.

Condition-averaged trajectories can be visualized in two-dimensional
latent space by plotting the first two latent components:

\begin{equation}
(z_1(t), z_2(t))
\end{equation}

Visualization typically includes

\begin{itemize}
\item condition-averaged trajectories
\item single-trial trajectories
\item markers indicating event onset
\item arrows indicating trajectory direction
\end{itemize}

Such plots illustrate how neural population states evolve through
latent space in response to experimental events.

\section{3D State-Space Visualization}

When the latent state space has more than two dimensions,
three-dimensional trajectory visualization can provide additional
insight.

A three-dimensional neural state space is defined by the first
three latent components:

\begin{equation}
(z_1(t), z_2(t), z_3(t))
\end{equation}

Neural trajectories can then be visualized as curves within this
three-dimensional space.

Key visualization elements include

\begin{itemize}
\item trajectories for different experimental conditions
\item color-coded time progression
\item trajectory start and end markers
\item rotation of the state-space view to reveal geometric structure
\end{itemize}

Three-dimensional visualization can reveal rotational dynamics,
trajectory divergence, and manifold structure that may not be
apparent in lower-dimensional projections.

\chapter{Computational Implementation}

The SEEG analysis framework described in this report is implemented
as a modular computational pipeline in MATLAB. The pipeline integrates
custom analysis functions with the FieldTrip toolbox, which provides
extensive functionality for electrophysiological data analysis.

The implementation is designed to support reproducible analysis of
large-scale SEEG datasets while allowing flexible extension for new
analysis modules.

\section{Pipeline Architecture}

The computational pipeline follows a modular architecture in which
each stage of the analysis is implemented as a separate processing
module.

The overall workflow is organized as follows:

\begin{enumerate}
\item Data loading and preprocessing
\item Event-locked trial extraction
\item Trial-level quality control
\item ERP analysis
\item Time--frequency analysis
\item Statistical inference
\item Multichannel population representation
\item Dimensionality reduction
\item Latent trajectory analysis
\item Dynamical systems analysis
\item Neural decoding
\item Visualization and reporting
\end{enumerate}

Each stage receives structured input data and produces output that
can be passed to downstream analysis modules.

Data are stored using structured MATLAB objects containing

\begin{itemize}
\item neural signals
\item trial metadata
\item channel labels
\item event timestamps
\item analysis parameters
\end{itemize}

This structure ensures consistent data flow across the entire
analysis pipeline.

\section{Software Design}

The software implementation emphasizes modularity and separation
of analysis components.

Each analysis step is implemented as a function that performs a
specific task. For example:

\begin{itemize}
\item \texttt{seeg\_preprocess.m} performs filtering and referencing
\item \texttt{seeg\_epoch.m} extracts event-locked trials
\item \texttt{seeg\_tf.m} computes time--frequency representations
\item \texttt{seeg\_latent.m} performs dimensionality reduction
\item \texttt{seeg\_trajectory.m} constructs neural trajectories
\item \texttt{seeg\_decode.m} performs neural decoding
\end{itemize}

This modular structure allows individual analysis components
to be tested and modified independently.

Integration with the FieldTrip toolbox enables efficient
implementation of commonly used signal processing functions,
including

\begin{itemize}
\item spectral analysis
\item wavelet transforms
\item statistical testing
\item sensor-level visualization
\end{itemize}

FieldTrip data structures are used where appropriate to maintain
compatibility with existing electrophysiological analysis tools.

\section{Reproducibility}

Reproducible analysis is essential for reliable scientific results.

To ensure reproducibility, the analysis pipeline incorporates
several design principles.

First, all analysis parameters are stored in configuration
structures that accompany each dataset. These configuration
files specify parameters such as

\begin{itemize}
\item filtering frequencies
\item time--frequency analysis parameters
\item dimensionality reduction settings
\item statistical thresholds
\end{itemize}

Second, intermediate results from each processing stage are
saved to disk. This enables reanalysis of later stages without
recomputing earlier processing steps.

Third, the entire analysis workflow can be executed using a
single master script that sequentially calls each analysis
module. This script ensures that all datasets are processed
using identical analysis parameters.

Finally, version control systems such as Git can be used to
track changes to the analysis code and maintain a history
of modifications to the computational pipeline.

Together, these design principles ensure that SEEG analyses
can be reproduced, validated, and extended in future studies.

\chapter{Complete MATLAB Pipeline Pseudocode}

This chapter summarizes the full SEEG computational analysis workflow
in MATLAB-style pseudocode. The goal is to provide a unified view of
the analysis pipeline from continuous SEEG recordings to neural
population dynamics, statistical inference, and visualization.

\section{Overall Pipeline}

The complete pipeline consists of the following stages:

\begin{enumerate}
\item Continuous signal preprocessing
\item Event-locked trial extraction
\item Trial-level quality control
\item ERP estimation
\item Time--frequency analysis
\item Oscillatory and phase-based feature extraction
\item Statistical inference
\item Multichannel population representation
\item Dimensionality reduction
\item Latent trajectory construction
\item Trajectory geometry and dynamical systems analysis
\item Neural decoding
\item Uncertainty quantification
\item Visualization and reporting
\end{enumerate}

\section{Main Analysis Routine}

\begin{algorithm}[H]
\caption{Complete SEEG analysis pipeline}
\begin{algorithmic}[1]
\Require Continuous SEEG data $X$, event table $E$, analysis configuration $\Theta$
\Ensure Analysis results structure $R$

\State Load continuous SEEG recording
\State Load channel metadata and event table
\State Initialize results structure $R$

\State $X_{clean} \gets$ \Call{PreprocessSignals}{$X, \Theta$}
\State $\mathcal{T} \gets$ \Call{ExtractEventLockedTrials}{$X_{clean}, E, \Theta$}

\State $\mathcal{Q} \gets$ \Call{ComputeTrialQualityMetrics}{$\mathcal{T}, \Theta$}
\State $(\mathcal{T}_{good}, \mathcal{W}) \gets$ \Call{ClassifyAndWeightTrials}{$\mathcal{T}, \mathcal{Q}, \Theta$}

\State $ERP \gets$ \Call{ComputeRobustERP}{$\mathcal{T}_{good}, \mathcal{W}, \Theta$}
\State $\mathcal{T}_{ind} \gets$ \Call{SubtractERP}{$\mathcal{T}_{good}, ERP$}

\State $TF_{total} \gets$ \Call{ComputeTimeFrequency}{$\mathcal{T}_{good}, \Theta$}
\State $TF_{ind} \gets$ \Call{ComputeTimeFrequency}{$\mathcal{T}_{ind}, \Theta$}
\State $ITPC \gets$ \Call{ComputeITPC}{$\mathcal{T}_{good}, \Theta$}

\State $F \gets$ \Call{ExtractNeuralFeatures}{$ERP, TF_{total}, TF_{ind}, ITPC, \Theta$}

\State $S_{ROI} \gets$ \Call{RunROIPermutationTests}{$F, \Theta$}
\State $S_{TF} \gets$ \Call{RunClusterTFStatistics}{$TF_{ind}, \Theta$}
\State $S_{onset} \gets$ \Call{EstimateResponseOnset}{$F, \Theta$}

\State $\mathcal{X}_{pop} \gets$ \Call{ConstructPopulationTensor}{$F, \Theta$}
\State $\mathcal{X}_{norm} \gets$ \Call{NormalizePopulationTensor}{$\mathcal{X}_{pop}, \Theta$}

\State $L \gets$ \Call{FitLatentModel}{$\mathcal{X}_{norm}, \Theta$}
\State $Z \gets$ \Call{ProjectToLatentSpace}{$\mathcal{X}_{norm}, L, \Theta$}
\State $Z_{smooth} \gets$ \Call{SmoothLatentTrajectories}{$Z, \Theta$}

\State $G \gets$ \Call{ComputeTrajectoryGeometry}{$Z_{smooth}, \Theta$}
\State $C \gets$ \Call{ComputeConditionSeparation}{$Z_{smooth}, \Theta$}
\State $D \gets$ \Call{RunDynamicalSystemsAnalysis}{$Z_{smooth}, \Theta$}

\State $Dec \gets$ \Call{RunNeuralDecoding}{$Z_{smooth}, \Theta$}
\State $U \gets$ \Call{BootstrapTrajectoryUncertainty}{$Z_{smooth}, \Theta$}

\State \Call{GenerateFigures}{$ERP, TF_{ind}, ITPC, Z_{smooth}, G, C, D, Dec, U$}

\State Store all outputs in $R$
\State \Return $R$
\end{algorithmic}
\end{algorithm}

\section{MATLAB-Oriented Function Layout}

A modular MATLAB implementation can be organized as follows:

\begin{verbatim}
run_seeg_pipeline.m
preprocess_signals.m
extract_event_locked_trials.m
compute_trial_quality.m
classify_trials.m
compute_robust_erp.m
compute_time_frequency.m
compute_itpc.m
extract_neural_features.m
run_roi_permutation_tests.m
run_cluster_tf_statistics.m
estimate_response_onset.m
construct_population_tensor.m
fit_latent_model.m
project_to_latent_space.m
compute_trajectory_geometry.m
compute_condition_separation.m
run_jpca_analysis.m
run_tangent_space_analysis.m
run_neural_decoding.m
bootstrap_trajectory_uncertainty.m
generate_figures.m
\end{verbatim}

\section{Master MATLAB Pseudocode}

\begin{verbatim}
function R = run_seeg_pipeline(X, E, Theta)

    X_clean   = preprocess_signals(X, Theta);
    Trials    = extract_event_locked_trials(X_clean, E, Theta);

    QC        = compute_trial_quality(Trials, Theta);
    [Trials, W] = classify_trials(Trials, QC, Theta);

    ERP       = compute_robust_erp(Trials, W, Theta);
    TrialsInd = subtract_erp(Trials, ERP);

    TFtotal   = compute_time_frequency(Trials, Theta);
    TFind     = compute_time_frequency(TrialsInd, Theta);
    ITPC      = compute_itpc(Trials, Theta);

    Feat      = extract_neural_features(ERP, TFtotal, TFind, ITPC, Theta);

    StatsROI  = run_roi_permutation_tests(Feat, Theta);
    StatsTF   = run_cluster_tf_statistics(TFind, Theta);
    StatsOn   = estimate_response_onset(Feat, Theta);

    Xpop      = construct_population_tensor(Feat, Theta);
    Xpop      = normalize_population_tensor(Xpop, Theta);

    Latent    = fit_latent_model(Xpop, Theta);
    Z         = project_to_latent_space(Xpop, Latent, Theta);
    Z         = smooth_latent_trajectories(Z, Theta);

    Geom      = compute_trajectory_geometry(Z, Theta);
    Sep       = compute_condition_separation(Z, Theta);
    Dyn       = run_dynamical_systems_analysis(Z, Theta);

    Dec       = run_neural_decoding(Z, Theta);
    Unc       = bootstrap_trajectory_uncertainty(Z, Theta);

    generate_figures(ERP, TFind, ITPC, Z, Geom, Sep, Dyn, Dec, Unc);

    R.ERP     = ERP;
    R.TF      = TFind;
    R.ITPC    = ITPC;
    R.Features= Feat;
    R.Stats   = struct('ROI',StatsROI,'TF',StatsTF,'Onset',StatsOn);
    R.Latent  = Z;
    R.Geometry= Geom;
    R.Separation = Sep;
    R.Dynamics   = Dyn;
    R.Decoding   = Dec;
    R.Uncertainty= Unc;

end
\end{verbatim}

\chapter{Algorithm Boxes for Each Analysis Stage}

This chapter summarizes each major analysis module as a standalone
algorithm box. These algorithmic descriptions are intended to clarify
the flow of computation and facilitate reproducible implementation.

\section{Signal Preprocessing}

\begin{algorithm}[H]
\caption{Signal preprocessing}
\begin{algorithmic}[1]
\Require Continuous SEEG data $X$
\Ensure Cleaned signal $X_{clean}$

\State Apply high-pass filter
\State Apply low-pass filter
\State Remove line noise with notch filters
\State Apply offline re-referencing
\State Detect bad channels using variance and noise metrics
\State Mark or remove artifact-contaminated segments
\State \Return $X_{clean}$
\end{algorithmic}
\end{algorithm}

\section{Event-Locked Signal Extraction}

\begin{algorithm}[H]
\caption{Event-locked trial extraction}
\begin{algorithmic}[1]
\Require Cleaned signal $X_{clean}$, event table $E$
\Ensure Trial tensor $\mathcal{T}$

\For{each event $e_k \in E$}
    \State Extract time window $[-T_{pre}, T_{post}]$
    \State Align extracted window to event onset
    \State Store trial data and metadata
\EndFor
\State \Return $\mathcal{T}$
\end{algorithmic}
\end{algorithm}

\section{Trial-Level Quality Control}

\begin{algorithm}[H]
\caption{Trial quality control}
\begin{algorithmic}[1]
\Require Trial tensor $\mathcal{T}$
\Ensure Trial quality labels and weights

\For{each trial}
    \State Compute prestimulus RMS amplitude
    \State Compute prestimulus variance
    \State Compute peak-to-peak amplitude
    \State Compute kurtosis
    \State Compute line-noise ratio
    \State Form composite quality score
    \State Assign trial to green, yellow, or red category
\EndFor
\State Assign analysis weights based on category
\State \Return quality labels and weights
\end{algorithmic}
\end{algorithm}

\section{ERP Analysis}

\begin{algorithm}[H]
\caption{Robust ERP estimation}
\begin{algorithmic}[1]
\Require Good trials $\mathcal{T}_{good}$ and weights $\mathcal{W}$
\Ensure ERP waveform $ERP(t)$

\For{each channel}
    \State Compute weighted average across trials
    \State Optionally reject high-leverage outliers
\EndFor
\State \Return ERP waveforms
\end{algorithmic}
\end{algorithm}

\section{Time--Frequency Analysis}

\begin{algorithm}[H]
\caption{Morlet wavelet time--frequency analysis}
\begin{algorithmic}[1]
\Require Trial signals $\mathcal{T}$
\Ensure Spectral power $P(t,f)$

\For{each trial}
    \For{each frequency $f$}
        \State Construct Morlet wavelet $\psi(t,f)$
        \State Convolve signal with wavelet
        \State Compute complex coefficient $Z(t,f)$
        \State Compute power $P(t,f)=|Z(t,f)|^2$
    \EndFor
\EndFor
\State Normalize power relative to baseline
\State \Return time--frequency representation
\end{algorithmic}
\end{algorithm}

\section{Phase-Based Synchronization}

\begin{algorithm}[H]
\caption{Inter-trial phase coherence}
\begin{algorithmic}[1]
\Require Complex wavelet coefficients $Z_k(t,f)$
\Ensure $ITPC(t,f)$

\For{each time-frequency point}
    \State Extract phase $\phi_k(t,f)$ for all trials
    \State Compute complex unit vectors $e^{i\phi_k(t,f)}$
    \State Average vectors across trials
    \State Take magnitude to obtain $ITPC(t,f)$
\EndFor
\State \Return ITPC map
\end{algorithmic}
\end{algorithm}

\section{ROI-Based Permutation Statistics}

\begin{algorithm}[H]
\caption{ROI permutation test}
\begin{algorithmic}[1]
\Require ROI feature values across trials and conditions
\Ensure Empirical p-value

\State Compute observed condition difference
\For{$b=1$ to $B$ permutations}
    \State Randomly shuffle condition labels
    \State Recompute ROI difference
\EndFor
\State Compare observed difference to null distribution
\State \Return p-value
\end{algorithmic}
\end{algorithm}

\section{Cluster-Based Time--Frequency Statistics}

\begin{algorithm}[H]
\caption{Cluster-based time--frequency permutation test}
\begin{algorithmic}[1]
\Require Time--frequency maps for two conditions
\Ensure Significant time--frequency clusters

\State Compute pointwise test statistic over time and frequency
\State Threshold statistic map
\State Form spatiotemporal clusters
\State Compute cluster mass for each cluster
\For{$b=1$ to $B$ permutations}
    \State Shuffle trial labels
    \State Recompute cluster masses
    \State Store maximum cluster mass
\EndFor
\State Compare observed cluster masses to permutation null
\State \Return significant clusters
\end{algorithmic}
\end{algorithm}

\section{Response Onset Detection}

\begin{algorithm}[H]
\caption{Permutation-based onset estimation}
\begin{algorithmic}[1]
\Require Time-resolved response statistic
\Ensure Response onset time $t_{onset}$

\State Compute observed timecourse statistic
\For{$b=1$ to $B$ permutations}
    \State Shuffle labels or sign-flip trials
    \State Recompute timecourse statistic
    \State Store maximum statistic across time
\EndFor
\State Determine corrected threshold
\State Find earliest timepoint above threshold
\State \Return $t_{onset}$
\end{algorithmic}
\end{algorithm}

\section{Multichannel Population Representation}

\begin{algorithm}[H]
\caption{Population tensor construction}
\begin{algorithmic}[1]
\Require Channel-wise features across trials and time
\Ensure Population tensor $\mathcal{X}$

\For{each trial}
    \For{each channel}
        \State Extract selected neural feature
        \State Insert feature into tensor
    \EndFor
\EndFor
\State Arrange data as trial $\times$ channel $\times$ time
\State \Return $\mathcal{X}$
\end{algorithmic}
\end{algorithm}

\section{Dimensionality Reduction}

\begin{algorithm}[H]
\caption{Latent state estimation}
\begin{algorithmic}[1]
\Require Population tensor $\mathcal{X}$
\Ensure Latent states $z(t)$

\State Reshape tensor into population state matrix
\State Fit dimensionality reduction model (PCA / FA / GPFA)
\State Project population activity into latent state space
\State \Return latent trajectories
\end{algorithmic}
\end{algorithm}

\section{Latent Neural Trajectories}

\begin{algorithm}[H]
\caption{Latent trajectory construction}
\begin{algorithmic}[1]
\Require Latent states across time and trials
\Ensure Single-trial and condition-averaged trajectories

\For{each trial}
    \State Construct sequence of latent states over time
\EndFor
\For{each condition}
    \State Average trajectories across trials
\EndFor
\State \Return latent trajectories
\end{algorithmic}
\end{algorithm}

\section{Trajectory Geometry}

\begin{algorithm}[H]
\caption{Trajectory geometry metrics}
\begin{algorithmic}[1]
\Require Latent trajectories $z(t)$
\Ensure Velocity, curvature, dispersion, and path length

\State Compute temporal derivatives of trajectory
\State Estimate velocity magnitude
\State Estimate curvature from first and second derivatives
\State Compute path length as cumulative displacement
\State Compute across-trial dispersion
\State \Return geometric metrics
\end{algorithmic}
\end{algorithm}

\section{Condition Separation}

\begin{algorithm}[H]
\caption{Condition separation in latent space}
\begin{algorithmic}[1]
\Require Condition-specific trajectories
\Ensure Time-resolved separation metric

\State Compute centroid trajectory for each condition
\For{each timepoint}
    \State Compute inter-condition distance
    \State Compute within-condition dispersion
    \State Form separation statistic
\EndFor
\State Assess significance using permutation testing
\State \Return separation timecourse
\end{algorithmic}
\end{algorithm}

\section{Rotational Dynamics}

\begin{algorithm}[H]
\caption{jPCA-based rotational dynamics analysis}
\begin{algorithmic}[1]
\Require Latent trajectory $z(t)$
\Ensure Rotational components and rotational plane

\State Compute temporal derivative $\dot{z}(t)$
\State Fit linear model $\dot{z}(t)=Mz(t)$
\State Constrain $M$ to skew-symmetric form
\State Compute eigenvectors of rotational dynamics matrix
\State Project latent states into rotational plane
\State \Return rotational trajectories
\end{algorithmic}
\end{algorithm}

\section{Tangent Space Dynamics}

\begin{algorithm}[H]
\caption{Tangent-space analysis}
\begin{algorithmic}[1]
\Require Latent trajectories for multiple conditions
\Ensure Tangent vectors and principal angle statistics

\State Estimate tangent vectors by finite differences
\State Build tangent subspace for each condition
\State Compute principal angles between subspaces
\State Optionally test significance by permutation
\State \Return tangent-space similarity metrics
\end{algorithmic}
\end{algorithm}

\section{Geometry of Neural Manifolds}

\begin{algorithm}[H]
\caption{Neural manifold geometry}
\begin{algorithmic}[1]
\Require Latent trajectories
\Ensure Manifold geometry descriptors

\State Compute covariance of latent states
\State Estimate state-space occupancy
\State Compute recurrence matrix
\State Estimate subspace overlap across conditions
\State \Return manifold geometry measures
\end{algorithmic}
\end{algorithm}

\section{Neural Decoding}

\begin{algorithm}[H]
\caption{Latent-state decoding}
\begin{algorithmic}[1]
\Require Latent states $z^{(k)}(t)$ and labels $y^{(k)}$
\Ensure Decoding accuracy and statistical significance

\For{each timepoint or analysis window}
    \State Split trials into training and test sets
    \State Train classifier on training set
    \State Evaluate classifier on test set
    \State Store decoding accuracy
\EndFor
\State Generate null accuracy distribution by label permutation
\State Estimate p-values and decoding onset
\State \Return decoding results
\end{algorithmic}
\end{algorithm}

\section{Uncertainty Quantification}

\begin{algorithm}[H]
\caption{Bootstrap uncertainty estimation}
\begin{algorithmic}[1]
\Require Single-trial trajectories
\Ensure Confidence intervals and confidence tubes

\For{$b=1$ to $B$ bootstrap iterations}
    \State Resample trials with replacement
    \State Recompute condition-averaged trajectory
    \State Recompute trajectory statistics
\EndFor
\State Compute percentile-based confidence intervals
\State Estimate covariance of bootstrap trajectories
\State Construct confidence tubes around mean trajectory
\State \Return uncertainty estimates
\end{algorithmic}
\end{algorithm}

\section{Visualization and Reporting}

\begin{algorithm}[H]
\caption{Visualization of neural dynamics}
\begin{algorithmic}[1]
\Require Analysis outputs from all stages
\Ensure Summary figures and report panels

\State Plot ERP waveforms with confidence intervals
\State Plot time--frequency power maps
\State Plot latent trajectories in 2D and 3D
\State Plot rotational planes and decoding curves
\State Plot confidence tubes and separation metrics
\State Export figures for report generation
\end{algorithmic}
\end{algorithm}

\chapter{Results}

\section{Overview}

Results were obtained from the derivative output directory
\texttt{/Volumes/OHDD/DATA/epilepsy/derivatives/seegring/sub-EP01AN96M1047}.
Two complementary result summaries were available. First, the batch
condition-discrimination report generated on 2026-03-24 summarized
population-level analyses across 9 cognitive tasks and 11 runs.
Second, a separate channel-significance summary quantified
channel-by-time cluster significance and decoding-based channel
importance for a subset of 7 analyzable task runs.

These two summaries support a consistent interpretation. At the
population latent-space level, no task showed statistically significant
condition separation after cluster-corrected permutation testing. In
contrast, at the channel level, several tasks exhibited substantial
numbers of channels with significant time clusters and robust decoding
performance, indicating that condition information was present in the
recordings but was not recovered as a statistically significant
population-level latent separation under the current analysis settings.

\section{Population-Level Condition Separation}

The cross-task report in
\texttt{reports/20260324\_171402/report.md} analyzed 9 cognitive tasks
spanning 11 runs, including lexical decision, shape control, sentence
noun, sentence grammar, salience pain, balloon watching, viseme
perception, and viseme generation. Statistical significance was assessed
with cluster-based permutation tests using 1000 permutations and
$\alpha=0.05$.

The main result of that report was negative: no task showed
statistically significant condition-dependent separation at the
population level. The corresponding summary table
(\texttt{stats\_summary.csv}) reported \texttt{Sig = NO} for every
processed task, with peak separation, minimum p-value, peak decoding,
decoding onset, effective dimensionality, and jPCA summary fields left
as undefined values. The report therefore does not support any claim of
significant condition separation in latent space, significant
condition-dependent temporal windows, or significant cross-condition
population decoding at the whole-task summary level.

\section{Channel-Level Significant Results}

A separate analysis stored in
\texttt{channel\_significance\_summary.mat} yielded a different picture
at the channel level. This analysis re-ran permutation statistics and
decoding for 7 task runs and computed the number of channels with
significant channel-by-time clusters. All 7 analyzed tasks showed at
least some cluster-corrected channel-level significance, although none
showed region-of-interest significance after false-discovery-rate
correction.

The strongest channel-level effects were observed for
\texttt{task-lexicaldecision\_run-01}, which showed 23 channels with
significant time clusters and a peak decoding accuracy of 92.0\%.
\texttt{task-shapecontrol\_run-01} showed a similarly strong effect,
with 22 significant channels and a peak decoding accuracy of 99.0\%.
Intermediate effects were observed for
\texttt{task-sentencegrammar\_run-01} and
\texttt{task-viseme\_run-01}, each with 9 significant channels, and for
\texttt{task-saliencepain\_run-01} and
\texttt{task-sentencenoun\_run-01}, each with 7 significant channels.
\texttt{task-balloonwatching\_run-01} showed 5 significant channels and
a nominal peak decoding accuracy of 100.0\%, although this result should
be interpreted cautiously given the very small sample size of 10 trials.

Across these 7 channel-level analyses, the count of ROI-level FDR
significant channels was zero in every task. Thus, the significant
effects reported here should be interpreted specifically as
channel-by-time cluster-corrected findings, not as ROI-level
FDR-significant contacts. The effects were spatially sparse and
time-localized rather than strong enough to survive ROI-level
multiple-comparison correction.

\section{Decoding Results}

The channel-level decoding analysis indicated that binary tasks carried
considerable condition information despite the absence of significant
population-level latent separation in the batch report. Peak decoding
accuracy reached 99.0\% in shape control, 92.0\% in lexical decision,
75.0\% in salience pain, 68.3\% in sentence grammar, and 65.0\% in
sentence noun. Balloon watching reached 100.0\%, but this estimate was
derived from only 10 trials and is therefore likely to be unstable.

By contrast, the multi-class viseme task showed much weaker decoding,
with a peak accuracy of 16.4\% across 14 classes. Although this is above
the nominal chance level of approximately 7.1\%, it was accompanied by
only modest channel-level clustering and was not reflected in a
significant cross-task population-level result.

\section{Spatial Distribution of Informative Channels}

The most repeatedly prominent channels across tasks were concentrated in
a relatively stable subset of contacts. Among the top-20 channels ranked
by decoding contribution, \texttt{PI6}, \texttt{PI7}, \texttt{PI5},
\texttt{PI4}, \texttt{Ia15}, and \texttt{PI3} each appeared in 6 of the
7 analyzed tasks. Additional frequently recurring channels included
\texttt{Ia14} in 5 tasks and \texttt{PI8}, \texttt{PI9}, \texttt{C14},
\texttt{C13}, \texttt{C12}, \texttt{Ia16}, and \texttt{PI2} in 4 tasks.

This recurrence suggests that a stable subset of electrodes carried much
of the discriminative signal across paradigms. The recurrent involvement
of neighboring \texttt{PI}, \texttt{Ia}, and \texttt{C} contacts is more
consistent with localized informative subspaces than with a broadly
distributed whole-population separation effect.

\section{Interpretation}

Taken together, the current outputs support a split conclusion. On the
one hand, the official batch condition-discrimination report provides no
evidence for statistically significant condition separation at the
population latent-space level after cluster-corrected testing. On the
other hand, the dedicated channel-significance analysis shows that
multiple tasks do contain statistically significant channel-by-time
effects and practically strong decoding, especially in lexical decision
and shape control.

The most conservative interpretation is therefore that condition-related
information is detectable in localized channels and in decoding space,
but that this information was not sufficiently strong, stable, or
distributed to yield significant whole-population latent separation in
the current batch summary. This distinction should be preserved in the
manuscript: significant channel-level effects are supported by the saved
outputs, but those channel-level effects do not survive the stricter
ROI-level FDR criterion, whereas significant population-level condition
separation is also not supported by the saved batch report.


\chapter{Discussion}

This chapter discusses the implications of population-level
SEEG analysis for understanding neural dynamics in the human brain.

\chapter{Conclusion}

The proposed computational framework integrates classical
electrophysiological analysis with modern neural population
dynamics approaches, enabling comprehensive investigation of
multichannel SEEG activity.



\end{document}
